% !TeX root = ../bd-screen.tex

\chapter{前言}

这是一本关于模态逻辑的导论性教材。我在卡尔加里大学教授“哲学~579.2（模态逻辑）”课程时，将其作为主要教材。它基于\href{https://openlogicproject.org}{开放逻辑项目}（Open Logic Project）的材料。

正文假定读者熟悉初等集合论以及（命题）逻辑的基础知识。这部分内容是我所授的另一门课程“逻辑~II”的先修内容。该课程的教材\href{https://slc.openlogicproject.org}{\emph{Sets, Logic, Computation}}（《集合、逻辑、计算》）也基于开放逻辑项目，因此亦可免费获取。不过，所需的这些背景材料已作为附录收录于本书中。我每次教授这门课程时，都会布置这些附录作为背景阅读。

\Cref{nml:::part}最初部分地基于 Aldo Antonelli（阿尔多·安东内利）关于“基本模态逻辑的经典对应理论”的讲义，这些讲义是他在 2015 年不幸早逝前贡献给开放逻辑项目的。我对这些笔记进行了大量修订和扩充，例如，关于框架可定义性和分析表的材料就是新增的。